% !TEX root = ../notes_template.tex
\chapter*{Tentative Schedule}
\addcontentsline{toc}{chapter}{Schedule}

\begin{center}
\begin{tabular}{|c|c|p{8cm}|}
\hline
\textbf{Date} & \textbf{Week} & \textbf{Topic} \\
\hline
June 2 & Week 1 & Movement \& Muscle Testing \\
\hline
June 9 & Week 2 & Intro to Reflexes, Pulses \& RUSI \\
\hline
June 16 & Week 3 & Shoulder Complex \\
\hline
June 23 & Week 4 & Elbow, Forearm, Wrist \& Hand \\
\hline
June 30 & Week 5 & Pelvis \& Hip \\
\hline
July 7 & Week 6 & Knee \\
\hline
July 14 & Week 7 & Ankle \& Foot \\
\hline
July 21 & Week 8 & Lumbopelvic \\
\hline
July 28 & Week 9 & Cervical \& Quarter Screens \\
\hline
Aug 4 & Week 10 & Clinical Skills Checklist \\
\hline
\end{tabular}
\end{center}

\bigskip
\noindent
\textit{Note:} Despite the regional organization of the laboratory, which is progressing anatomically from the shoulder complex to the cervical spine, its foundation of clinical reasoning remains grounded in a regional interdependence framework. The regional structure ensures that the lab content aligns directly with the concurrent lecture material and the Osteology \& Arthrology Lab, supporting consistency throughout the summer. However, you will be continually encouraged to think beyond the local region being tested or imaged. Dysfunction in one area is rarely isolated; mobility or motor control deficits often manifest through compensations between joints and systems. Thus, while we examine one anatomical region at a time, our diagnostic and intervention strategies must account for whole-body function and multiregional interactions.

